A rigorous partial resolution is available. Fix \(K\geq2\) and a nonzero zero-sum contrast \(c\).

If \(|S_c|=2\), write the two active coefficients as \(\alpha\) and \(-\alpha\), where \(\alpha=L_c/2\). Independent assignment with probability \(1/2\) on each active arm, the estimator \(\alpha f_{n,2}^\star\) obtained from the published scalar posterior-mean rule, and the internally lifted complete-schedule least-favorable prior form an exact finite-\(n\) minimax saddle against arbitrary dependent designs and arbitrary possibly biased estimators. Hence
\[
 q^\star=(1/2,1/2)\quad\hbox{on }S_c,
 \qquad
 \rho_n(c)=C_0(c)\rho_n((1,-1)),
\]
and
\[
 d_n(c)=C_0(c)C_A n^{-4/3}+o(n^{-4/3}).
\]

For arbitrary support, the independent \(q^\star\) design and the explicit clipped-shrinkage rule from the universal second-order rate theorem satisfy, for every sufficiently large \(n\),
\[
 \max_{z\in\mathcal T_K^n}
 R_n(\mathcal D^\star,\widehat\tau_n^{\mathrm{sh}};z)
 \leq C_0(c)\{n^{-1}-\kappa_c n^{-4/3}\},
 \qquad \kappa_c>0,
\]
and the unrestricted embedded two-arm converse gives
\[
 0<C_0(c)\kappa_c
 \leq\liminf_{n\to\infty}n^{4/3}d_n(c)
 \leq\limsup_{n\to\infty}n^{4/3}d_n(c)
 \leq C_0(c)C_A.
\]
Thus the universal second-order exponent is exactly \(4/3\).

When \(|S_c|\geq3\), the subsequential-limit set of \(n^{4/3}d_n(c)\) is a nonempty compact subset of \((0,\infty)\). For rational \(c\), taking \(M=n\) in the exact-rational finite program gives a computable orbit procedure and a rational-prior converse whose risk gap is at most \(C_0(c)/(4n^2)=o(n^{-4/3})\); their normalized improvements have exactly the same subsequential-limit set and convergence behavior as \(n^{4/3}d_n(c)\). For every real \(c\), one can choose rational zero-sum contrasts \(c^{(n)}\) with \(L_{c-c^{(n)}}/2=o(n^{-5/6})\) and again take \(M=n\), obtaining transferred computable upper and lower certificates whose gap is \(o(n^{-4/3})\) and whose normalized improvement endpoints have the same subsequential-limit set and convergence behavior as \(n^{4/3}d_n(c)\). These finite programs retain every response type and require no exposure-positivity, allocation-support, rank, or overlap restriction on the unrestricted minimax design.

For \(K=3\) and \(c=c^\dagger=(1,-1/2,-1/2)\), one has \(C_0(c^\dagger)=1\), and the exact three-arm certificates satisfy, for every \(n,M\geq1\),
\[
 B_n(\nu_{n,M})\leq\rho_n(c^\dagger)\leq U_{n,M}
 \leq B_n(\nu_{n,M})+\frac{1}{4M^2}.
\]
The exact-rational calculations for \(n\leq5\) are finite-sample diagnostics for this sandwich only.

For \(|S_c|\geq3\), these conclusions do not determine whether \(n^{4/3}d_n(c)\) converges, its exact constant or exact subsequential-limit set, second-order optimality of \(q^\star\), convergence of finite-program optimizers or least-favorable priors, an analytic multivariate feedback rule, or a face--corner Hamilton--Jacobi--Bellman, spectral, separable, or other limiting operator. Those questions remain unresolved.